% ============================================================
% CAPITOLO 40 - LA REGOLA DELL'OTTAVA COMPLETA (CROMATICA)
% ============================================================
\chapter{La Regola dell'Ottava completa (cromatica)}
\label{chap:regola_ottava_cromatica}

La Regola dell'Ottava diatonica (Capitolo \ref{chap:regola_ottava_completa}) lascia al VI grado discendente un'armonizzazione semplice, \sigla{IV}{6}{}. Le tabelle complete prevedono qui anche accordi cromatici, che questo capitolo aggiunge.

\section{Alterazioni nella Regola}
\label{sec:alterazioni_regola}

\begin{regola}[Varianti cromatiche della Regola dell'Ottava]
	\label{reg:regola_cromatica}
	Nella scala \textbf{discendente} il VI grado al basso può ricevere un accordo cromatico:
	\begin{itemize}
		\item \textbf{modo maggiore:} \sigla{II}{4}{3} con la terza alzata (cifra $\sharp$6 4 3), cioè la \textbf{dominante secondaria di V} (in Do maggiore Re-Fa$\sharp$-La-Do con il La al basso). Non è una sesta eccedente: quella in maggiore si forma sul La$\flat$ (VI grado abbassato) ed è trattata nel Capitolo \ref{chap:seste_eccedenti};
		\item \textbf{modo minore:} la \textbf{sesta eccedente} (cifra $\sharp$6), con il Fa al basso (La minore), il Re$\sharp$ nella voce superiore, che risolve espandendosi verso l'ottava Mi.
	\end{itemize}
\end{regola}

\begin{esempio}[Regola dell'Ottava in Sol maggiore: scala discendente con la dominante secondaria di V]
	\label{ex:regola_cromatica_magg}
	Sul VI grado discendente (Mi) la cifra $\sharp$6 4 3 indica la dominante secondaria di V (La-Do$\sharp$-Mi-Sol con il Mi al basso), che risolve su V.
	\musica{es_cap14_01}
\end{esempio}

\begin{esempio}[Regola dell'Ottava in La minore: scala discendente con la sesta eccedente]
	\label{ex:regola_cromatica_min}
	Sul VI grado discendente (Fa) la cifra $\sharp$6 indica la sesta italiana, con il Re$\sharp$ che sale verso il Mi mentre il basso scende al Mi.
	\musica{es_cap14_03}
\end{esempio}

\begin{table}[h]
	\centering
	\caption{Regola dell'Ottava: VI grado discendente, variante diatonica e cromatica}
	\label{tab:regola_ottava_cromatica}
	\small
	\begin{tabular}{lccc}
		\toprule
		\textbf{Modo} & \textbf{Basso} & \textbf{Diatonica} & \textbf{Cromatica} \\
		\midrule
		Maggiore & La & 6 (\sigla{IV}{6}{}) & $\sharp$6 4 3 (V/V) \\
		Minore & Fa & 6 (\sigla{iv}{6}{}) & $\sharp$6 (sesta eccedente) \\
		\bottomrule
	\end{tabular}
\end{table}
